% ===== PRÉAMBULE CANONIQUE — à coller VERBATIM en tête de CHAQUE .tex =====
\documentclass[11pt,a4paper]{article}
\usepackage[utf8]{inputenc}\usepackage[T1]{fontenc}\usepackage{lmodern}
\usepackage[french]{babel}
\usepackage{amsmath,amssymb,mathtools}\usepackage{icomma}
\usepackage[version=4]{mhchem}          % \ce{...} : équations & formules
\usepackage{siunitx}\sisetup{locale=FR} % \qty{}{}, \si{}, \ang{}
\usepackage{chemfig}                    % \chemfig{...} : structures développées
\usepackage{xcolor}\usepackage{colortbl}
\usepackage{tikz}\usepackage{pgfplots}\pgfplotsset{compat=1.18}
\usepackage[most]{tcolorbox}\usepackage{enumitem}\usepackage{array,booktabs}
\usepackage[a4paper,margin=2.2cm]{geometry}
\usetikzlibrary{arrows.meta,calc,angles,quotes,positioning,patterns,backgrounds,shapes.geometric}
\definecolor{primary}{RGB}{222,49,99}\definecolor{primarydark}{RGB}{150,22,55}
\definecolor{defcol}{RGB}{0,114,178}\definecolor{methcol}{RGB}{52,92,148}
\definecolor{excol}{RGB}{0,135,90}\definecolor{attncol}{RGB}{200,40,55}
\tcbset{boxbase/.style={enhanced,breakable,boxrule=0.9pt,arc=2.5pt,
  left=3mm,right=3mm,top=2.3mm,bottom=2.3mm,fonttitle=\bfseries,coltitle=white}}
% --- boîtes TUTORIEL (noms EXACTS, ne pas renommer) ---
\newtcolorbox{cle}[1][]{boxbase,colback=defcol!6,colframe=defcol,
  title={\textsc{À retenir}\ifx\relax#1\relax\relax\else\textmd{ : \itshape #1}\fi}}
\newtcolorbox{methode}[1][]{boxbase,colback=methcol!7,colframe=methcol,sharp corners,
  title={\textsc{Méthode}\ifx\relax#1\relax\relax\else\textmd{ --- \itshape #1}\fi}}
\newtcolorbox{exemplebox}[1][]{boxbase,colback=excol!5,colframe=excol,
  title={\textsc{Exemple}\ifx\relax#1\relax\relax\else\textmd{ : \itshape #1}\fi}}
\newtcolorbox{piege}[1][]{boxbase,colback=attncol!6,colframe=attncol,
  title={\textsc{Piège}\ifx\relax#1\relax\relax\else\textmd{ : \itshape #1}\fi}}
\newtcolorbox{remarque}[1][]{boxbase,colback=black!4,colframe=black!45,
  title={\textsc{Remarque}\ifx\relax#1\relax\relax\else\textmd{ : \itshape #1}\fi}}
% --- boîtes EXERCICE / CORRIGÉ (noms EXACTS, auto-comptées) ---
\newtcolorbox[auto counter]{exo}[1][]{boxbase,colback=primary!4,colframe=primary,
  title={\textsc{Exercice}~\thetcbcounter\ifx\relax#1\relax\relax\else\textmd{ --- \itshape #1}\fi}}
\newtcolorbox[auto counter]{corrige}[1][]{boxbase,colback=excol!4,colframe=excol,
  title={\textsc{Corrigé}~\thetcbcounter\ifx\relax#1\relax\relax\else\textmd{ --- \itshape #1}\fi}}
\newcommand{\R}{\mathbb{R}}\newcommand{\N}{\mathbb{N}}\newcommand{\Z}{\mathbb{Z}}\newcommand{\ds}{\displaystyle}
% (un \usepackage{fancyhdr}+en-tête est possible ; il est ignoré côté web)
% =========================================================================
\begin{document}
\sisetup{per-mode=symbol}
\begin{center}\small\itshape
Aux TPN : $V_{\mathrm{m}}=\qty{22.4}{\litre\per\mole}$. Rendement \qty{100}{\percent} sauf mention contraire.
\end{center}

\begin{exo}[doser les réactifs d'une synthèse]
On veut préparer \qty{0.300}{\mole} de sulfate d'aluminium par
$\ce{2 Al(OH)3 + 3 H2SO4 -> Al2(SO4)3 + 6 H2O}$. On donne
$M(\ce{Al(OH)3})=\qty{78}{\gram\per\mole}$ et $M(\ce{H2SO4})=\qty{98}{\gram\per\mole}$.
\begin{enumerate}[label=\alph*)]
  \item Quelle \textbf{masse} de \ce{Al(OH)3} faut-il engager ?
  \item Quelle \textbf{masse} de \ce{H2SO4} faut-il engager ?
  \item Quelle quantité d'eau se forme ?
\end{enumerate}
\end{exo}

\begin{exo}[tableau d'avancement avec un réactif en excès]
On brûle \qty{0.2}{\mole} d'éthanol dans un \emph{excès} de dioxygène :
$\ce{C2H5OH + 3 O2 -> 2 CO2 + 3 H2O}$.
\begin{enumerate}[label=\alph*)]
  \item Exprime $n(\ce{C2H5OH})$, $n(\ce{CO2})$ et $n(\ce{H2O})$ en fonction de $\xi$.
  \item L'oxygène étant en excès, détermine l'avancement maximal $\xi_{\max}$.
  \item Quel volume de \ce{CO2} (aux TPN) obtient-on ?
\end{enumerate}
\end{exo}

\begin{exo}[combustion du kérosène : la chaîne complète]
Le kérosène \ce{C14H30} ($M=\qty{198}{\gram\per\mole}$, $\rho=\qty{0.763}{\gram\per\milli\litre}$)
brûle selon $\ce{2 C14H30 + 43 O2 -> 28 CO2 + 30 H2O}$. On en brûle \qty{100}{\milli\litre}.
\begin{enumerate}[label=\alph*)]
  \item Quelle masse, puis quelle quantité de matière de kérosène brûle-t-on ?
  \item Quelle quantité de \ce{CO2} se forme-t-il ?
  \item Quel volume de \ce{CO2} cela représente-t-il aux TPN ?
\end{enumerate}
\end{exo}

\begin{exo}[dosage par titrage]
On dose une solution de dichromate \ce{Cr2O7^2-} de concentration inconnue par une solution d'ions
\ce{Fe^2+} à \qty{0.60}{\mole\per\litre}, selon
$\ce{Cr2O7^2- + 6 Fe^2+ + 14 H+ -> 2 Cr^3+ + 6 Fe^3+ + 7 H2O}$.
On titre \qty{10}{\milli\litre} de dichromate ; l'équivalence est atteinte pour \qty{12}{\milli\litre}
d'ions \ce{Fe^2+}.
\begin{enumerate}[label=\alph*)]
  \item Écris la relation entre $n(\ce{Cr2O7^2-})$ et $n(\ce{Fe^2+})$ à l'équivalence.
  \item Calcule la quantité d'ions \ce{Fe^2+} versée à l'équivalence.
  \item En déduire la concentration $[\ce{Cr2O7^2-}]$.
\end{enumerate}
\end{exo}

\end{document}
